% ============================================================
% TikZ EXAMPLES — Ví dụ minh họa cách dùng tikz-math.tex
% ============================================================
% File này KHÔNG được \input vào main.tex
% Dùng để tham khảo khi cần vẽ cho bài toán cụ thể
% ============================================================

% ==========================================
% VÍ DỤ 1: BẢNG BIẾN THIÊN — HÀM BẬC 3
% ==========================================
%
% Hàm y = x^3 - 3x^2 + 2, CĐ tại x=0, CT tại x=2
%
% \bbtBacBa{0}{2}{2}{-2}{-\infty}{+\infty}
%
% Tham số:
%   #1 = x cực đại     (0)
%   #2 = x cực tiểu     (2)
%   #3 = f(x1) cực đại  (2)
%   #4 = f(x2) cực tiểu (-2)
%   #5 = lim x→-∞       (-\infty)
%   #6 = lim x→+∞       (+\infty)


% ==========================================
% VÍ DỤ 2: BẢNG BIẾN THIÊN — HÀM BẬC 4
% ==========================================
%
% Hàm y = x^4 - 2x^2 + 1
%   CT tại x=-1 và x=1 → f = 0
%   CĐ tại x=0 → f = 1
%
% \bbtBacBon{-1}{0}{1}{0}{1}{0}{+\infty}{+\infty}
%
% Tham số:
%   #1 = x cực tiểu trái  (-1)
%   #2 = x cực đại         (0)
%   #3 = x cực tiểu phải   (1)
%   #4 = f(x1)             (0)
%   #5 = f(x2)             (1)
%   #6 = f(x3)             (0)
%   #7 = lim x→-∞          (+\infty)
%   #8 = lim x→+∞          (+\infty)


% ==========================================
% VÍ DỤ 3: BBT TỰ DO (custom)
% ==========================================
%
% Dùng khi hàm có nhiều hơn 3 cực trị hoặc dạng đặc biệt
%
% \begin{bangBT}{12}{-5}
%   % Tự vẽ theo ý muốn, dùng các style: bbt, bbt header, bbt arrow, bbt val, bbt sign
%   \draw[bbt] (0,-1) -- (12,-1);   % đường kẻ ngang
%   \draw[bbt] (1.5,0) -- (1.5,-5); % đường kẻ dọc
%   \node[bbt header] at (0.75, -0.5) {$x$};
%   % ... thêm nội dung tùy ý
% \end{bangBT}


% ==========================================
% VÍ DỤ 4: ĐỒ THỊ HÀM SỐ (pgfplots)
% ==========================================
%
% \begin{center}
% \begin{tikzpicture}
%   \begin{axis}[
%     forum graph,             % hoặc forum graph small
%     xmin=-3, xmax=5,
%     ymin=-5, ymax=5,
%     xtick={-2,-1,...,4},
%     ytick={-4,-2,...,4},
%   ]
%     % Vẽ hàm số
%     \addplot[TikzCurve1, domain=-2:4] {x^3 - 3*x^2 + 2};
%
%     % Tiệm cận đứng
%     \tiemcanDung{2}{-5}{5}
%
%     % Tiệm cận ngang
%     \tiemcanNgang{1}{-3}{5}
%
%     % Điểm đặc biệt (đặc, fill)
%     \diemDB{0}{2}{$(0;2)$}{above right}
%
%     % Điểm rỗng (hở)
%     \diemHo{2}{-2}{$(2;-2)$}{below}
%
%     % Đường thẳng y = m
%     \addplot[TikzCurve2, dashed, domain=-2:4] {1};
%
%     % Legend
%     \legend{$y=f(x)$, $y=m$}
%   \end{axis}
% \end{tikzpicture}
% \end{center}


% ==========================================
% VÍ DỤ 5: ĐỒ THỊ HÀM PHÂN THỨC
% ==========================================
%
% y = (2x+1)/(x-1), TCĐ: x=1, TCN: y=2
%
% \begin{center}
% \begin{tikzpicture}
%   \begin{axis}[
%     forum graph small,
%     xmin=-5, xmax=7,
%     ymin=-6, ymax=10,
%   ]
%     % Nhánh trái
%     \addplot[TikzCurve1, domain=-5:0.85] {(2*x+1)/(x-1)};
%     % Nhánh phải
%     \addplot[TikzCurve1, domain=1.15:7] {(2*x+1)/(x-1)};
%
%     % Tiệm cận
%     \tiemcanDung{1}{-6}{10}
%     \tiemcanNgang{2}{-5}{7}
%
%     % Giao trục
%     \diemDB{-0.5}{0}{$(-\frac{1}{2};0)$}{above left}
%     \diemDB{0}{-1}{$(0;-1)$}{below right}
%   \end{axis}
% \end{tikzpicture}
% \end{center}


% ==========================================
% VÍ DỤ 6: VÙNG TÍCH PHÂN
% ==========================================
%
% Tô vùng giữa đồ thị f(x) = x^2 và trục Ox, từ x=0 đến x=2
%
% \begin{center}
% \begin{tikzpicture}
%   \begin{axis}[
%     forum graph small,
%     xmin=-1, xmax=3,
%     ymin=-0.5, ymax=5,
%   ]
%     % Tô vùng
%     \addplot[TikzFill, opacity=0.3, domain=0:2] {x^2} \closedcycle;
%
%     % Vẽ đường cong phía trên
%     \addplot[TikzCurve1, domain=-0.5:2.5] {x^2};
%
%     \diemDB{0}{0}{$O$}{below left}
%     \diemDB{2}{0}{$2$}{below}
%   \end{axis}
% \end{tikzpicture}
% \end{center}


% ==========================================
% VÍ DỤ 7: HÌNH HỌC KHÔNG GIAN — HÌNH CHÓP
% ==========================================
%
% Hình chóp S.ABCD
%
% \begin{center}
% \begin{tikzpicture}[scale=1.2]
%   % Đáy ABCD (hình bình hành, phối cảnh)
%   \dinh{A}{(-2, 0)}{below left}
%   \dinh{B}{(1.5, -0.5)}{below right}
%   \dinh{C}{(3, 1)}{right}
%   \dinh{D}{(-0.5, 1.5)}{above left}
%   \dinh{S}{(0.5, 4)}{above}
%
%   % Cạnh đáy
%   \draw[solid edge] (A) -- (B) -- (C);
%   \draw[hidden edge] (C) -- (D) -- (A);
%
%   % Cạnh bên
%   \draw[solid edge] (S) -- (A);
%   \draw[solid edge] (S) -- (B);
%   \draw[solid edge] (S) -- (C);
%   \draw[hidden edge] (S) -- (D);
%
%   % Chiều cao
%   \coordinate (H) at (0.5, 0.3); % chân đường cao
%   \draw[height line] (S) -- (H);
%   \node[font=\footnotesize, color=ForumRed] at (0.1, 2) {$h$};
%   \gocVuong{H}{A}{B}{0.25cm}
% \end{tikzpicture}
% \end{center}


% ==========================================
% VÍ DỤ 8: MẶT PHẲNG PHỨC
% ==========================================
%
% Biểu diễn z = 2 + 3i và |z - (1+i)| = 2
%
% \begin{center}
% \begin{tikzpicture}
%   \begin{axis}[
%     forum graph small,
%     axis equal,
%     xlabel={Re}, ylabel={Im},
%     xmin=-2, xmax=5,
%     ymin=-2, ymax=5,
%   ]
%     % Đường tròn |z-(1+i)|=2
%     \addplot[complex circle, domain=0:360, samples=100]
%       ({1 + 2*cos(x)}, {1 + 2*sin(x)});
%
%     % Vùng tô bên trong
%     \addplot[complex region, domain=0:360, samples=100]
%       ({1 + 2*cos(x)}, {1 + 2*sin(x)}) \closedcycle;
%
%     % Điểm z = 2+3i
%     \diemDB{2}{3}{$z=2+3i$}{above right}
%
%     % Tâm
%     \diemDB{1}{1}{$I(1;1)$}{below right}
%   \end{axis}
% \end{tikzpicture}
% \end{center}


% ==========================================
% VÍ DỤ 9: CÂY XÁC SUẤT
% ==========================================
%
% P(A) = 0.3, P(B|A) = 0.5, P(B|Ac) = 0.2
%
% \begin{center}
% \begin{tikzpicture}[
%   level 1/.style={sibling distance=4cm, level distance=2.5cm},
%   level 2/.style={sibling distance=2cm, level distance=2.5cm},
% ]
%   \node[prob node] {Start}
%     child {
%       node[prob node] {$A$}
%         child {
%           node[prob result] {$A \cap B$}
%           edge from parent node[prob label, left] {$0{,}5$}
%         }
%         child {
%           node[prob result] {$A \cap \bar{B}$}
%           edge from parent node[prob label, right] {$0{,}5$}
%         }
%       edge from parent node[prob label, left] {$0{,}3$}
%     }
%     child {
%       node[prob node] {$\bar{A}$}
%         child {
%           node[prob result] {$\bar{A} \cap B$}
%           edge from parent node[prob label, left] {$0{,}2$}
%         }
%         child {
%           node[prob result] {$\bar{A} \cap \bar{B}$}
%           edge from parent node[prob label, right] {$0{,}8$}
%         }
%       edge from parent node[prob label, right] {$0{,}7$}
%     };
% \end{tikzpicture}
% \end{center}
